\section{\iflanguage{english}{Conclusion}{Fazit}}\label{sec:conclusion}

In dieser Arbeit wurde untersucht, ob und wie gut SATZilla – ein auf maschinellem Lernen basierendes 
Algorithmus-Portfolio – für die Auswahl zwischen verschiedenen Konfigurationen des modernen 
CDCL-SAT-Solvers Kissat eingesetzt werden kann. Die Experimente umfassten drei Hauptphasen: 
(1) Validierung der Pipeline mit künstlichen ("Fake")-Solvern, (2) Training und Evaluation auf sechs 
realen Kissat-Varianten mit den ursprünglichen Benchmarks sowie (3) Analyse der Auswirkungen von 
Vorverarbeitung auf die Vorhersagen und Laufzeiten.

Erstens bestätigt die erfolgreiche Validierung mit Fake-Solvern, dass die SATZilla-Pipeline – von 
der Feature-Extraktion über das Training bis zur Vorhersage – korrekt implementiert ist. Die hohe 
Genauigkeit und die plausible Features beweisen, dass das Verfahren lernen kann.

Zweitens zeigt die Anwendung auf sechs echte Kissat-Varianten, dass SATZilla auch bei komplexen, 
realen Solvern nutzbare Vorhersagen liefert. Mit einer Top-3-Genauigkeit von 85 \% und einem 
durchschnittlichen Speedup von 1,47 gegenüber dem Standard-Solver ist SATZilla nicht nur besser 
als ein Zufallsverfahren, sondern auch praktisch nützlich für die portfoliogestützte SAT-Lösung. 
Die Analyse der Features offenbart zudem, welche strukturellen Eigenschaften einer 
Instanz für die Leistungsunterschiede der Kissat-Varianten verantwortlich sind – ein Erkenntnisgewinn, 
der über die reine Solver-Auswahl hinausgeht.

Drittens belegt die Untersuchung der Vorverarbeitung, dass SATZillas Vorhersagen nicht 
vorverarbeitungsinvariant sind. Die Änderungen der Solver-Empfehlungen sind eine direkte Folge der 
tatsächlichen strukturellen Veränderungen der Formeln, nicht willkürlich. Die Dosis-Wirkungs-Beziehung 
(mehr Vorverarbeitung → mehr Änderungen) validiert die Sensitivität des Feature-Sets. Für die 
praktische Anwendung bedeutet dies, dass SATZilla in diesem Fall Vorhersagen macht die Zutreffender 
sein könnten.

Viertens zeigt der Laufzeitvergleich, dass Vorverarbeitung kein Allheilmittel ist: Sie beschleunigt 
zwar etwa die Hälfte der Instanzen, kann in Einzelfällen aber zu extremen Verlangsamungen führen. 
Ein signifikanter Vorteil von 100.000 gegenüber 10.000 Vereinfachungsschritten konnte nicht 
festgestellt werden.

Abschließend möchte ich sagen das mir die Arbeit sehr gefallen hat und ich viel über ein neues Thema 
lernen konnte. Besonder spannend fand ich die Arbeit an der Forschung von Vorhersagen und die anschließende
Auswertung der Daten. 

\subsubsection{Technische Daten}

Die Arbeit wurde von meinem heim PC, Windows 11, mit WSl gemacht, mit diesem konnte ich mich mit einer ssh verbindung 
am PC der Technischen Fakultät anmelden, dies war von großem vorteil da ich breits im meinem 
Bachelorprojekt mehr RAM für meine brechnungen bruachte. Tmux erwies sich als nützlich Programme über 
einen längeren Zeitraum laufen zu lassen.

Das Programmiern der Skripte wurde mit Python 3.9 geschrieben mit folgenden Bibliotheken: numpy, pandas, 
scikit-learn (random forest), matplotlib (zum erstellen der Diagramme und Visualierungen), seaborn, joblib, plotly, 
tqdm, requests. Ausgeführt wurden die Skripte mit Conda.




